\documentclass[10pt,a4paper]{amsart}
\usepackage[margin=19mm]{geometry}
\usepackage{amsmath,amssymb,mathtools}
\usepackage[scr=rsfs,cal=euler]{mathalfa}
\usepackage{xfrac}
\usepackage[unicode,hidelinks,bookmarks=false]{hyperref}
\newcommand{\ie}{i.e.\ }
\newcommand{\cf}{cf.\ }
\newcommand{\resp}{respectively\ }
\newcommand{\Subsection}{\subsection}
\providecommand{\nobreakdash}{\nobreak\hbox{-}}
\setlength{\emergencystretch}{2em}
\multlinegap0pt
\usepackage{bm}
\usepackage{bbm}
\usepackage{dsfont}
\usepackage{xspace}
\usepackage{cancel}
\usepackage{mathtools}
\usepackage{amsthm}
\makeatletter
\@ifundefined{lemma}{\newtheorem{lemma}{Lemma}}{}
\makeatother
\makeatletter
\newcommand{\dd}{{\mathrm d}}
\expandafter\ifx\csname ecn\endcsname\relax
\newenvironment{ecn}{\begin{equation}}{\end{equation}}
\fi
\expandafter\ifx\csname eqnarr\endcsname\relax
\newenvironment{eqnarr}{\begin{IEEEeqnarray}{rCl}}{\end{IEEEeqnarray}\ignorespacesafterend}
\fi
\renewcommand{\eqref}[1]{\hyperref[#1]{(\ref*{#1})}}
\expandafter\ifx\csname esn\endcsname\relax
\newenvironment{esn}{\begin{equation*}}{\end{equation*}}
\fi
\expandafter\ifx\csname listeo\endcsname\relax
\newenvironment{listeo}{\begin{list}{\alph{enumi})}{\usecounter{enumi}}}{\end{list}}
\fi
\newcommand*{\norm}[1]{\lVert #1 \rVert}
\expandafter\ifx\csname numerai\endcsname\relax
\newenvironment{numerai}{\begin{list}{\roman{enumi})}{\usecounter{enumi}}}{\end{list}}
\fi
\expandafter\ifx\csname shadedbox\endcsname\relax
\newenvironment{shadedbox}[1][]%
        {
        %\setlength{\fboxsep}{-\fboxrule}
        %\footnotesize\normalfont\ttfamily
          \raggedright
        \setlength{\rightmargin}{\leftmargin}
        \setlength{\itemsep}{-12pt}
        \setlength{\parsep}{20pt}
        \begin{lrbox}{\@tempboxa}%
        \begin{minipage}{\linewidth-2\fboxsep}
        }%
        {
        \end{minipage}%
        \end{lrbox}%
        \fcolorbox{black}{shade}{\usebox{\@tempboxa}}\newline
        }%
\fi
\makeatother
\title[The replicator coalescent]{The replicator coalescent}
\author{A.E. Kyprianou, L. Pe\~naloza, T. Rogers}
\date{Source: arXiv:2207.00998v3 (CC BY 4.0)}
\begin{document}
\maketitle

\noindent\emph{Source and licence.} The replicator coalescent. Version arXiv:2207.00998v3 (CC BY 4.0), distributed under the Creative Commons Attribution 4.0 International licence, as recorded in the arXiv licence metadata for this exact version and shown on the abstract page https://arxiv.org/abs/2207.00998v3. Full rights notice: licenses/arxiv-2207.00998-CC-BY-4.0.txt.\par\smallskip

\noindent\textit{Editorial scope.} Counted unit: the Lemma whose source label is \texttt{cor1} in the article (source0.tex lines 1253--1258, source label \texttt{cor1}), delivered together with the article's own complete original proof of it (source0.tex lines 1260--1292, one \texttt{proof} environment of 33 lines). No mathematics is written, repaired or re-derived: statement and proof are the article's own text line for line. Required context retained uncounted: comingdownfrominfinity (lines 795--808); rhoupper (lines 981--990); dt (lines 1148--1151). Every definition, display and label of those lines is retained unchanged. Omitted from this package and not redistributed: 1--794, 809--980, 991--1147, 1152--1252, 1259--1259, 1293--2039. Nothing inside a retained excerpt is omitted or altered: every retained body is delivered exactly as the article's lines are written, so no omission receipt of any kind accompanies this package. Citations: the retained excerpts carry no citation command at all, so no bibliography entry of the article is transcribed and no reference is redistributed. Reference adaptation: none was needed and none was made; every reference inside the delivered unit resolves inside it, so no unresolved reference or citation is produced. Printed slips kept literal: no printed word, symbol, hypothesis or number of the counted statement or of its proof was altered. Layout and class: the article's own class is \texttt{article} with packages including amsmath, amsthm, enumitem, thmtools, amsfonts, amssymb, color, graphicx, tcolorbox, wrapfig, caption, hyperref, tikz, pgfplots; the wrapper is the lane's pinned \texttt{amsart} editorial wrapper at 10pt on A4, which restates the theorem-like environments the retained text needs and adds the article's own macro definitions that the retained text needs (\texttt{\textbackslash dd}, \texttt{\textbackslash eqref}, \texttt{\textbackslash norm}), with extra wrapper packages (bm, bbm, dsfont, xspace, cancel, mathtools). The delivered numbering is the unit's own. Assets: no retained span contains a figure environment, a graphics-inclusion command, a PDF-page inclusion, a table or a TikZ picture; no image file is copied into this package, the package declares no asset dependency and no image file is redistributed with it. Licence: version arXiv:2207.00998v3 of this article is distributed under the Creative Commons Attribution 4.0 International licence, as recorded in the arXiv licence metadata for this exact version and shown on the abstract page https://arxiv.org/abs/2207.00998v3. A full rights notice is delivered at licenses/arxiv-2207.00998-CC-BY-4.0.txt. Characters: every retained excerpt is pure ASCII, so the delivered engine, XeLaTeX with its default Latin Modern text font, reports no missing character.
\begin{lemma}\label{comingdownfrominfinity}
Without any further requirement on $C_{i,j}$, there exists a Kingman coalescent death chain $\nu^-$ such that $\sigma\geq \nu^-$ under $\mathbb{P}_{\bm\eta}$, for each ${\bm \eta} = (\arg({\bm n}(0)), \sigma(0)^{-1})\in \mathcal{S}^k\times \mathbb{N}^{-1}$. 
Conversely, if $C_{i,j}>0$ for all $i,j$, there exists a Kingman coalescent death chain $\nu^+$ such that $\sigma\leq \nu^+$ under $\mathbb{P}_{\bm\eta}$, for each ${\bm \eta} = (\arg({\bm n}(0)), \sigma(0)^{-1})\in \mathcal{S}^k\times \mathbb{N}^{-1}$. 

In particular, under the  assumption that $C_{i,j}>0$ for all $i,j$,
the replicator coalescent comes down from infinity  in the sense that, for all $\varepsilon>0$, 
%\[
%\lim_{N\to\infty}\mathbb{P}_{\bm\eta^N}(0<\gamma_m<\infty)=1, \qquad \forall\, m\in \mathbb{N},
%\]
%and
\[
\lim_{m\to\infty}\lim_{N\to\infty}\mathbb{P}_{\bm\eta^N}(\gamma_m<\varepsilon) = 1.
\]
\end{lemma}

\begin{equation}
\rho(\bm{n})
<
\sum_{i=1}^k\frac{1}{2}Cn_i\left(\sum_{j\neq i}n_{j} +  n_i-1\right) = 
\sum_{i=1}^k\frac{1}{2}Cn_i(\norm{\bm n}_1-1) = {C} {\norm{\bm n}_1 \choose 2}.
%
%
% C{\norm{\bm n}_1 \choose 2} - \sum_{i=1}^k \frac{(C_{i,i}-\underline{C})}{2}n_i> C{\norm{\bm n}_1 \choose 2} - \frac{(\overline{C}-\underline{C})}{2}\norm{\bm n}_1,
\label{rhoupper}
\end{equation}

\begin{equation}
\int_0^{\tau(t)} \sigma(s)\dd s = t \,\,\,\text{ and hence } \,\,\, \sigma(\tau(t))\dd \tau(t) = \dd t %\,\,\, \text{ for }\,\,\,  t<\int_0^{\gamma_1}\sigma(s)\dd s=:\zeta.
\label{dt}.
\end{equation}

\begin{lemma}\label{cor1} Fix $t>0$. Suppose $(\bm \eta^N, N\geq 1)$ tends to $( \bm r_0, 0)$,  
\[
\tau^{-1}(t) =   \int^{t}_0 \sigma(s)\dd s \to \infty, \, \,  \tau(t)\to 0,\, \, \text{ and }\,\, |\tau(t)\wedge \gamma_1-\tau(t)|\to0,
\]
weakly as $N\to\infty$.
\end{lemma}

\begin{proof}
Recall from the proof of Lemma \ref{comingdownfrominfinity}, and specifically \eqref{rhoupper},  there is a death chain $(\nu(t), t\geq0)$, representing the number of blocks in a Kingman coalescent with merger rate $C$  such that, for any $\bm \eta\in\mathbb{N}^k_*$, on the same probability space, we can stochastically bound $\sigma(t)\geq \nu(t)$, $t\geq0$.


\medskip


We now note that, for any large $M>0$ there exists a constant $C>0$ (not necessarily the same as before) such that, for any $m$ sufficiently large,
\begin{align}
\lim_{N\to\infty}\mathbb{P}_{{\bm \eta}^N}\left(\int_0^{\gamma_m} \sigma(s)\dd s> M\right) 
&\geq \lim_{N\to\infty}\mathbb{P}_{{\bm \eta}^N}\left(\int_0^{\beta_m} \nu(s)\dd s> M\right) \notag\\
&= \Pr\left(\sum_{n = m+1}^\infty \frac{n}{C{n\choose 2}}\mathbf{e}^{(n)}_{1} >M\right)\notag\\
&=\Pr\left(\sum_{n = m+1}^\infty \frac{1}{n-1}\mathbf{e}^{(n)}_{1} >\frac{CM}{2}\right),
\label{eq1}
\end{align}
where $\beta_m = \inf\{t>0: \nu(t) = m\}$ and $({\bf e}^{n}_1, n\geq 1)$, is a sequence of iid unit-mean exponentially distributed random variables.
If we write $(N(t), t\geq 0)$ for a Poisson process with unit rate, then almost surely we have
\[
\sum_{n = m+1}^\infty \frac{1}{n-1}\mathbf{e}^{(n)}_{1} = \int_0^\infty \frac{1}{N(s)+m}\dd s 
=\int_0^\infty \frac{s}{N(s)+m}\frac{\dd s}{s} =\infty
\]
where  the final equality follows by the strong law of large numbers for Poisson processes. As such the right-hand side of \eqref{eq1} is equal to 1.

\medskip


Since $M$ and $m$ can be arbitrarily large, this shows the first claim as soon as we note that $\tau^{-1}(t) = \int_0^t \sigma(u)\dd u$, which  is an easy consequence of  \eqref{dt}. On the other hand, note that since 
$
\int_0^{\tau(t)} \sigma(s)\dd s = t,
$
when $\sigma(0)= {\bm \eta}^N \rightarrow\infty$, the above comparison with Kingman's coalescent shows that  almost surely that, since $\int_0^u\sigma(s)\dd s $ converges weakly to infinity for all $u>0$, then $\tau(t)$ converges weakly to 0. Indeed, if with positive probability $\tau(t)>\varepsilon$ in the limit as $N\to\infty$, then, on that event, 
$\int_0^{\tau(t)} \sigma(s)\dd s\geq \int_0^{\varepsilon} \sigma(s)\dd s$, which explodes in distribution. This in turn contradicts the definition of $\tau(t)$. This proves the second and third statements of the lemma.
\end{proof}

\end{document}
